% Zone „Das kennst du schon“ – aus bank/vektoren-und-rechenoperationen/zone.jsonl (zusammenbau v0.1)
\einheitenkopf[zone][Kennst du schon]{Das kennst du schon}
\setcounter{aufgabe}{0}

%% TODO Titel: Ich-kann-Satz fehlt in der Bank (Platzhalter: Kettenname) – Nr. 1
%% TODO Anweisung: ein Satz über der ersten Teilaufgabe fehlt in der Bank – Nr. 1
\begin{aufgabe}{Verschiebungen im ebenen Koordinatensystem ausführen und Punkte ablesen}
\begin{teile}
\teil $P(2 | 1)$ um 3 nach rechts und 4 nach oben: Bildpunkt? P'( \leerfeld | \leerfeld )
\teil $Q(-3 | 2)$ um 5 nach rechts und 6 nach unten: Bildpunkt? Q'( \leerfeld | \leerfeld )
\end{teile}
\end{aufgabe}

%% TODO Titel: Ich-kann-Satz fehlt in der Bank (Platzhalter: Kettenname) – Nr. 2
%% TODO Anweisung: ein Satz über der ersten Teilaufgabe fehlt in der Bank – Nr. 2
\begin{aufgabe}{Satz des Pythagoras und Wurzelgleichungen (quadrieren, Lösungen sieben)}
\begin{teile}
\teil Katheten 3 cm und 4 cm: Hypotenuse? \leerfeld[cm]
\teil Katheten 1,5 m und 2 m: Hypotenuse? \leerfeld[m]
\end{teile}
\end{aufgabe}

%% TODO Titel: Ich-kann-Satz fehlt in der Bank (Platzhalter: Kettenname) – Nr. 3
%% TODO Anweisung: ein Satz über der ersten Teilaufgabe fehlt in der Bank – Nr. 3
\begin{aufgabe}{Punkte im räumlichen Koordinatensystem lesen, Schrägbilder von Quadern und Prismen deuten}
\begin{teile}
\teil Koordinaten der Ecke, die dem Ursprung gegenüberliegt?

\begin{ksys3} \rquader{0,0,0}{4}{3}{2} \end{ksys3}
( \leerfeld | \leerfeld | \leerfeld )
\teil Lies die Koordinaten von $Q$ ab.

\begin{ksys3} \rpunkt{3}{-2}{2}{Q} \end{ksys3}
Q( \leerfeld | \leerfeld | \leerfeld )
\end{teile}
\end{aufgabe}

%% TODO Titel: Ich-kann-Satz fehlt in der Bank (Platzhalter: Kettenname) – Nr. 4
%% TODO Anweisung: ein Satz über der ersten Teilaufgabe fehlt in der Bank – Nr. 4
\begin{aufgabe}{Terme mit einem Buchstaben umformen und lineare Gleichungen lösen}
\begin{teile}
\teil $3r = 12$: $r$? r = \leerfeld
\teil $2{,}5t + 1{,}5t = 12$: $t$? t = \leerfeld
\end{teile}
\end{aufgabe}

%% TODO Titel: Ich-kann-Satz fehlt in der Bank (Platzhalter: Kettenname) – Nr. 5
%% TODO Anweisung: ein Satz über der ersten Teilaufgabe fehlt in der Bank – Nr. 5
\begin{aufgabe}{Mit Anteilen und Verhältnissen rechnen (ein Verhältnis mit einem Faktor ansetzen)}
\begin{teile}
\teil Saft : Wasser = 1 : 3, zusammen 8 Liter: wie viel Saft? \leerfeld[l]
\teil Rot : Gelb : Blau = 2 : 3 : 5, zusammen 30 ml: jede Farbe? Rot \leerfeld[ml] , Gelb \leerfeld[ml] , Blau \leerfeld[ml]
\end{teile}
\end{aufgabe}

%% TODO Titel: Ich-kann-Satz fehlt in der Bank (Platzhalter: Kettenname) – Nr. 6
%% TODO Anweisung: ein Satz über der ersten Teilaufgabe fehlt in der Bank – Nr. 6
\begin{aufgabe}{Skalarprodukt in Koordinaten berechnen}
\begin{teile}
\teil $(1 | 2 | 3) \circ (2 | 1 | 1)$ = ?
\teil $(2 | -3 | 1) \circ (4 | 2 | -5)$ = ?
\end{teile}
\end{aufgabe}

%% TODO Titel: Ich-kann-Satz fehlt in der Bank (Platzhalter: Kettenname) – Nr. 7
%% TODO Anweisung: ein Satz über der ersten Teilaufgabe fehlt in der Bank – Nr. 7
\begin{aufgabe}{Satz des Pythagoras und Wurzelgleichungen (quadrieren, Lösungen sieben) – Fehler finden}
\begin{teile}
\teil Ich finde den Fehler: Tom löst $x^2 = 64$ und schreibt $x = 8$.
\end{teile}
\end{aufgabe}

%% TODO Titel: Ich-kann-Satz fehlt in der Bank (Platzhalter: Kettenname) – Nr. 8
%% TODO Anweisung: ein Satz über der ersten Teilaufgabe fehlt in der Bank – Nr. 8
\begin{aufgabe}{Satz des Pythagoras und Wurzelgleichungen (quadrieren, Lösungen sieben) – selbst rechnen}
\begin{teile}
\teil $(x - 1)^2 = 9$: alle Lösungen? x = \leerfeld oder x = \leerfeld
\end{teile}
\end{aufgabe}
